\addcontentsline{toc}{section}{Вывод}

\begin{center}
	\section*{Вывод}
\end{center}

В курсовой работе был рассмотрен протокол прикладного уровня DHCP, рассмотрены его положительные и отрицательные стороны, компоненты и основные принципы работы. Также рассмотрен метод его построения, который был воспроизведен на практике при помощи ISC-DHCP-Server на операционной системе Linux Fedora и с использованием виртуальной машины. Можно отметить, что создание такого локального сервера не занимает много времени, является достаточно легким и нетребовательно к техническим характеристикам компьютера, что является несомненно одним из самых главных преимуществ.